\documentclass[11pt]{article}
\usepackage{amsmath,amssymb,amsthm,mathtools}
\usepackage[T1]{fontenc}
\usepackage{lmodern}
\usepackage[margin=1in]{geometry}
\usepackage{microtype}
\usepackage{xcolor}
\usepackage{hyperref}
\usepackage{xurl}

\hypersetup{colorlinks=true,linkcolor=blue,citecolor=blue,urlcolor=blue}
\urlstyle{same}

\newtheorem{theorem}{Theorem}[section]
\newtheorem{proposition}[theorem]{Proposition}
\newtheorem{remark}[theorem]{Remark}

\title{Promoted-collar breaker:\\ complete local audit and finite active bursts}
\author{}
\date{March 25, 2026}

\begin{document}
\maketitle

\section{Promoted-collar breaker: complete local audit and finite active bursts}

Let
\[
\widehat{\mathfrak N}_m(1)
=
\{A_{1,4},\,C_{r,4},\,A_{1,3},\,C_{3,3},\,C_{r-1,3},\,A_{r,3},\,C_{r+1,3}\},
\qquad r=(m-1)/2,
\]
and for \(u\in\{r-1,r+1\}\) define the promoted-collar package
\[
\widehat{\mathfrak N}_m^{(u)}:=\widehat{\mathfrak N}_m(1)\cup\{A_{u,3}\}.
\]
Write
\[
R^{(u)}_{1,m}=h_1^m\big|_{\Sigma=0}.
\]
Recall
\[
\beta=\mathbf 1[c=m-2],
\qquad
\chi=\mathbf 1[d+\beta\equiv m-1],
\]
and the exact promoted-collar identity
\[
R^{(u)}_{1,m}=R^{\mathrm{wd3}}_{1,m}+\chi(e_2-e_0).
\]

\begin{proposition}[exact active locus]
\label{prop:promoted-active-locus}
One has
\[
\chi=1
\iff
\bigl(d=m-1,\ c\neq m-2\bigr)\ \text{or}\ \bigl(d=c=m-2\bigr).
\]
Hence the promoted-collar correction is supported exactly on
\[
\mathcal A_m
=
\{d=m-1,\ c\neq m-2\}\sqcup\{d=c=m-2\}.
\]
Outside \(\mathcal A_m\), one has
\[
R^{(u)}_{1,m}=R^{\mathrm{wd3}}_{1,m}=R^{\mathrm{wd1}}_{1,m}.
\]
\end{proposition}

\begin{proposition}[top hinge]
\label{prop:promoted-top-hinge}
Assume \(c=d=m-2\) and \(e\neq m-1\). Then
\[
R^{(u)}_{1,m}(a,b,m-2,m-2,e)
=
\begin{cases}
(5,\ *,\ 1,\ m-1,\ e), & a=2,\\[1mm]
(a+2,\ *,\ 2,\ m-1,\ e), & a\neq 2,
\end{cases}
\]
where \(*\) is determined by \(\Sigma=0\). In particular, every top-hinge state enters the top-collar interior after one block.
\end{proposition}

\begin{proof}
Here \(\beta=1\), \(\chi=1\), \(e\neq m-1\), so \(\varepsilon=0\), \(\alpha=1\), \(\gamma=\delta=0\), and
\[
\phi=\mathbf 1[a\equiv 2].
\]
Hence the promoted exact block map gives
\[
\Delta_0=2+\phi,\qquad \Delta_2=4-\phi,\qquad \Delta_3=1,\qquad \Delta_4=0,
\]
from which the formula follows immediately.
\end{proof}

\begin{proposition}[double hinge]
\label{prop:promoted-double-hinge}
Assume \(c=d=m-2\) and \(e=m-1\). Then
\[
R^{(u)}_{1,m}(a,b,m-2,m-2,m-1)
=
\begin{cases}
(5,\ *,\ 1,\ m-1,\ m-1), & a=3,\\[1mm]
(a+1,\ *,\ 2,\ m-1,\ m-1), & a\neq 3,
\end{cases}
\]
where \(*\) is determined by \(\Sigma=0\). In particular, every double-hinge state enters the double-top interior after one block.
\end{proposition}

\begin{proof}
Now \(\beta=1\), \(\chi=1\), \(e=m-1\), so \(\varepsilon=1\), \(\alpha=0\), \(\gamma=\delta=0\), and
\[
\phi=\mathbf 1[a\equiv 3].
\]
Hence
\[
\Delta_0=1+\phi,\qquad \Delta_2=4-\phi,\qquad \Delta_3=1,\qquad \Delta_4=0,
\]
which is exactly the stated formula.
\end{proof}

\begin{theorem}[finite promoted-collar bursts]
\label{thm:promoted-finite-bursts}
For every odd resonant modulus \(m\ge 15\), every maximal orbit segment of
\(R^{(u)}_{1,m}\) contained in
\[
\mathcal A_m
=
\{d=m-1,\ c\neq m-2\}\sqcup\{d=c=m-2\}
\]
is finite.

More precisely, a segment starting on the hinge line enters either the top-collar interior or the double-top interior after one block, and thereafter exits that interior after finitely many further blocks by the exact top-collar and double-top forcing theorems.
\end{theorem}

\begin{proof}
The support decomposition is given by Proposition~\ref{prop:promoted-active-locus}, and the one-step hinge entrances are Propositions~\ref{prop:promoted-top-hinge} and \ref{prop:promoted-double-hinge}. The previously established promoted-collar transducer theorems imply that uninterrupted top-collar and double-top episodes are finite. Combining these facts gives the claim.
\end{proof}

\begin{remark}[what remains]
The finite-burst theorem completes the local section-return audit: there is no hidden persistent \(\chi\)-active locus beyond the top interior, double-top interior, and hinge line. What remains is global: one still has to prove that inserting these finite bursts into the solved width-1 bulk produces a single lower cycle, and one still has to control the donor-color side of the final resonant package.
\end{remark}

\end{document}
